% !TEX root = ../main.tex
\chapter{Planteamiento del problema y justificación}

La histopatología permite caracterizar enfermedades mediante la observación de
la morfología microscópica de los tejidos. Su digitalización en forma de WSI
facilita conservar, consultar y distribuir las láminas, y proporciona una entrada
estandarizada para desarrollar herramientas de análisis computacional
~\cite{pallua_future_2020,song_artificial_2023}. Estas herramientas pueden apoyar
tareas repetitivas, cuantificar patrones y explorar biomarcadores morfológicos,
pero su adopción clínica exige distinguir el desempeño experimental de la
validación diagnóstica y considerar limitaciones operativas, regulatorias y de
generalización~\cite{verghese_computational_2023}.

Una dificultad central es la disponibilidad de etiquetas. Las anotaciones
histológicas detalladas requieren tiempo de especialistas y su costo aumenta con
el tamaño de la lámina, la complejidad de la estructura y la resolución espacial
solicitada. En un corpus anotado de colon y piel se observaron tiempos entre 45
segundos y más de 360 minutos por WSI, según la tarea
~\cite{lindman_annotations_2019}. En consecuencia, es común disponer de muchas
regiones de tejido que pueden utilizarse para aprender una representación, pero
de un número menor de regiones con etiquetas apropiadas para entrenar y evaluar
un clasificador.

Las redes neuronales profundas aprenden representaciones mediante composiciones
de transformaciones parametrizadas y no requieren que las características se
definan manualmente de antemano~\cite{goodfellow_deep_2016}. Para imágenes, las
CNN aprovechan la localidad y la compartición espacial de pesos: un mismo núcleo
se aplica en diferentes posiciones, lo que reduce el número de parámetros frente
a una capa densa de dimensiones comparables y produce equivarianza traslacional
en la operación convolucional. Esta propiedad no implica invariancia completa de
la red y puede verse afectada por bordes, muestreo, agrupamiento y pasos mayores
que uno~\cite{cohen_group_2016}.

Las rotaciones constituyen una simetría adicional relevante en histopatología,
pues una orientación distinta de una región no debería cambiar su contenido
diagnóstico. El aumento de datos mediante rotaciones expone al modelo a ejemplos
transformados, pero no garantiza una respuesta exactamente equivariante y añade
un costo que depende de cuántas transformaciones se evalúen. Las convoluciones
equivariantes incorporan la relación geométrica en la arquitectura mediante
compartición restringida de parámetros~\cite{veeling_rotation_2018}. Esta
restricción puede reducir la complejidad estadística bajo supuestos adecuados,
aunque su utilidad debe comprobarse en la tarea y distribución concretas
~\cite{elesedy_provably_2021}.

Por tal motivo, este trabajo estudia redes neuronales convolucionales
equivariantes como parte de un esquema semisupervisado. Primero se entrenaron,
sin consultar las etiquetas de las tareas posteriores, un VAE convencional y un
VAE con equivarianza continua a \(\mathrm{SO}(2)\). Después se congelaron ambos
codificadores y se comparó la utilidad de sus representaciones mediante
clasificación diagnóstica a nivel de WSI y segmentación de la WSI sobre una
cuadrícula cuya unidad es un parche de \(256\times256\) píxeles. Esta última
evaluación se restringió a regiones anotadas de alta pureza; las zonas sin
anotación no se trataron como una clase negativa.

Esto conduce a la siguiente pregunta de investigación:

¿En qué medida las representaciones aprendidas por un VAE con equivarianza
continua a \(\mathrm{SO}(2)\), comparadas con las de un VAE convolucional base,
resultan útiles para la clasificación a nivel de WSI y la segmentación de WSI a
resolución de parche cuando se dispone de anotaciones limitadas?
